\section{The grid}
\label{sec:grid}

\subsection{A recursive array of any dimension}

In Java the dimension of an array is part of its type: \code{int[][]} has two dimensions and
always will. The subject therefore asks for a data structure whose dimension $d$ is an input of
the program, defined recursively: an array of size $k_1 \times k_2 \times \dots \times k_n$
allocates $k_1$ arrays of size $k_2 \times \dots \times k_n$.

\code{DynamicArrayND<T>} is that structure. The definition has two cases, a recursive one and a
base one, and the class models them with two private subclasses instead of a flag and a pair of
half-used fields:

\begin{itemize}
    \item a \code{Branch} holds a list of arrays that have one dimension less;
    \item a \code{Leaf} is a one-dimensional array and holds the elements.
\end{itemize}

This is the Composite pattern. Client code only sees the abstract type and obtains an instance
from a static factory method, whose signature uses a variable number of arguments, as the subject
suggests:

\begin{lstlisting}
DynamicArrayND<GridCell> cells = DynamicArrayND.create(GridCell::new, 65, 65, 65);
GridCell center = cells.get(32, 32, 32);
\end{lstlisting}

The first argument computes the element stored at given coordinates; it is called once per
element, so an array is never observed half-filled. The construction follows the recursive
definition literally:

\begin{lstlisting}
private static <T> DynamicArrayND<T> allocate(
        Shape shape, int depth, int[] prefix,
        Function<? super Coordinates, ? extends T> initializer) {
    int length = shape.size(depth);
    if (depth == shape.dimension() - 1) {                      // base case: a row of elements
        Object[] elements = new Object[length];
        for (int index = 0; index < length; index++) {
            prefix[depth] = index;
            elements[index] = initializer.apply(Coordinates.of(prefix));
        }
        return new Leaf<>(elements);
    }
    List<DynamicArrayND<T>> children = new ArrayList<>(length);  // recursive case
    for (int index = 0; index < length; index++) {
        prefix[depth] = index;
        children.add(allocate(shape, depth + 1, prefix, initializer));
    }
    return new Branch<>(children);
}
\end{lstlisting}

Access is recursive in the same way. The public method \code{get(int... indices)} checks that
there is one index per dimension, then calls \code{get(indices, 0)}; a \code{Branch} picks the
child designated by \code{indices[depth]} and delegates to it with \code{depth + 1}; a
\code{Leaf} returns the element. An access costs $d$ steps and allocates nothing. The class is
generic and implements \code{Iterable<T>}, visiting elements in row-major order.

\begin{figure}[tbp]
    \centering
    \begin{tikzpicture}
        \node[class] (array) at (0, 0) {\atitle{DynamicArrayND<T>}
            \nodepart{second}+ \static{create}(initializer, sizes...)\\
            + get(indices...): T\\+ set(value, indices...)\\+ dimension(): int\\+ shape(): Shape\\
            + iterator(): Iterator<T>\\\# \textit{get(indices, depth): T}};
        \node[class] (leaf) at (-2.2, -3.5) {\ctitle{Leaf<T>}\nodepart{second}- elements: Object[]};
        \node[class] (branch) at (1.75, -3.5) {\ctitle{Branch<T>}
            \nodepart{second}- children: List<DynamicArrayND<T>>};

        \node[class=3] (grid) at (6.8, 0.25) {\ctitle{Grid}
            \nodepart{second}- shape: Shape\\- boundary: BoundaryCondition\\- cells: List<GridCell>
            \nodepart{third}+ cellAt(coordinates...): Cell\\+ cell(coordinates...): GridCell\\
            + setState(position, state)\\+ evolve(transition): int\\+ population(): int\\
            + snapshot(): int[]\\+ fingerprint(): long};
        \node[enum=3] (boundary) at (10.9, 0.95) {\ctitle[enumeration]{BoundaryCondition}
            \nodepart{second}FIXED\\PERIODIC
            \nodepart{third}\textasciitilde{} \textit{outside(...): Cell}};

        \node[iface] (cell) at (10.9, -2.5) {\ctitle[interface]{Cell}
            \nodepart{second}+ state(): int\\+ isOn(): boolean};
        \node[class=3] (gridcell) at (6.8, -4.6) {\ctitle{GridCell}
            \nodepart{second}- position: Coordinates\\- state: int\\- staged: int
            \nodepart{third}+ position(): Coordinates\\\textasciitilde{} stage(nextState)\\
            \textasciitilde{} commit(): boolean};
        \node[enum] (void) at (10.9, -5.0) {\ctitle[enumeration]{VoidCell}\nodepart{second}INSTANCE};

        \node[class] (slice) at (-1.2, -6.4) {\ctitle{Slice}
            \nodepart{second}+ \static{parse}(text): Slice\\+ extract(grid): Frame\\
            + locate(row, column): Coordinates\\+ moved(axis, coordinate): Slice};
        \node[class] (frame) at (3.6, -6.9) {\ctitle{Frame}
            \nodepart{second}+ rows, columns: int\\+ stateAt(row, column): int};

        \draw[extends] (leaf.north) -- ++(0, 0.45) -| ($(array.south)+(-0.4,0)$);
        \draw[extends] (branch.north) -- ++(0, 0.45) -| ($(array.south)+(0.4,0)$);
        \draw[composition] (branch.east) -- ++(0.3, 0) |- node[elabel, pos=0.25] {$k_1$ children}
        ($(array.east)+(0,-0.9)$);
        \draw[composition] ($(grid.west)+(0,0.4)$) -- node[elabel] {stores cells in}
        ($(array.east)+(0,0.65)$);
        \draw[association] ($(grid.east)+(0,0.7)$) -- (boundary.west);
        \draw[composition] (grid.south) -- node[elabel] {$k_1 \cdots k_d$} (gridcell.north);
        \draw[implements] (gridcell.east) -- ++(0.9, 0) |- (cell.west);
        \draw[implements] (void) -- (cell);
        \draw[uses] (boundary.south) -- node[elabel] {resolves to} (cell.north);
        \draw[uses] (slice.east |- frame.west) -- node[elabel, above] {produces} (frame.west);
    \end{tikzpicture}
    \caption{The \code{grid} package. \code{Shape} and \code{Coordinates}, two small immutable value
        types used throughout, are not drawn.}
    \label{fig:grid}
\end{figure}

\paragraph{Why not a flat array?} A single \code{int[]} indexed by
$i_1 k_2 \cdots k_d + \dots + i_d$ would be more compact and faster. The recursive structure was
kept because the subject asks for it, and it was isolated behind a class of its own: nothing else
in the program knows how cells are stored, so the representation could be swapped without
touching the rest. The measurements of \cref{sec:quality} show that it is fast enough.

\subsection{Cells, and what lies beyond the edge}

A \code{Grid} owns a \code{DynamicArrayND<GridCell>}. Each \code{GridCell} knows its position and
its state (an \code{int}, so that automata are not limited to two states).

Automata are defined on an infinite grid; the subject settles the finite case by stating that any
cell outside the grid is permanently off. Sprinkling bounds checks over every operator would have
been error-prone. Instead, the question ``which cell is at these coordinates?'' always has an
answer:

\begin{lstlisting}
public Cell cellAt(int... coordinates) {
    if (shape.contains(coordinates)) {
        return storage.get(coordinates);
    }
    return boundary.outside(this, coordinates);
}
\end{lstlisting}

Rules see cells through the two-method interface \code{Cell}. Outside the grid, the
\code{FIXED} boundary answers with \code{VoidCell.INSTANCE}, a cell whose state is 0 forever: a
Null Object, implemented as a single-element enumeration, which is the idiomatic Java singleton.
The boundary condition itself is a Strategy written as an enumeration with one method body per
constant; the second constant, \code{PERIODIC}, wraps coordinates around and turns the grid into
a torus. No rule, operator or neighborhood contains a bounds check.

\subsection{Simultaneous evolution}

``All the cells of a cellular automaton evolve simultaneously.'' Updating cells in place, one
after the other, would let a cell read the \emph{new} state of a neighbor updated before it and
silently compute a different automaton. The grid owns the protocol that prevents this:

\begin{lstlisting}
public int evolve(ToIntFunction<? super GridCell> transition) {
    if (cells.size() >= PARALLEL_THRESHOLD) {
        cells.parallelStream().forEach(cell -> cell.stage(transition.applyAsInt(cell)));
    } else {
        for (GridCell cell : cells) {
            cell.stage(transition.applyAsInt(cell));
        }
    }
    int changed = 0;
    for (GridCell cell : cells) {
        if (cell.commit()) {
            changed++;
        }
    }
    return changed;
}
\end{lstlisting}

In the first phase every cell computes its next state from the untouched current generation and
\emph{stages} it in a private field; in the second, all cells \emph{commit}. \code{stage} and
\code{commit} are package-private: only the grid can call them, so the invariant cannot be broken
from outside.

The first phase only reads shared state and writes one private field per cell. It is therefore
safe to run on several threads, which the grid does for 8\,192 cells or more, using a parallel
stream over the flat list of cells built once at construction. The number of changed cells,
returned for free by the second phase, is what detects convergence.

A unit test pins the property down: with the transition ``copy your left neighbor'', a row
\code{1 0 0 0 0} must become \code{0 1 0 0 0}, not \code{1 1 1 1 1}.

\subsection{Slices}

A grid of three dimensions or more is displayed through a plane. \code{Slice} parses the
Matlab-like notation of the subject, \code{(37, :, :)}, where a colon marks a displayed axis, and
\code{extract(grid)} returns the selected plane as a \code{Frame}, an immutable two-dimensional
picture of states. The same class handles the degenerate cases, \code{(:, :)} for a 2D grid and
\code{(:)} for a 1D one, so the display code is identical in every dimension.
